\begin{table}[tbp]
\caption{Six-factor alphas before and after the betting-against-beta factor.}
\label{tab:alpha}
\centering\small
\setlength{\tabcolsep}{5pt}
\begin{tabular}{lrrrrrr}
\toprule
Portfolio & $\alpha$ (\%) & $t$ & $\alpha$+BAB (\%) & $t$ & $\beta_{\rm BAB}$ & $t$\\
\midrule
USMV, the fund & -0.17 & -0.10 & -1.96 & -1.45 & 0.309 & 12.2\\
USMV replica, structural & -0.07 & -0.09 & -0.54 & -0.71 & 0.099 & 10.3\\
USMV replica, sample & 0.58 & 0.66 & 0.16 & 0.19 & 0.113 & 8.8\\
USMV replica, Ledoit--Wolf & 0.91 & 0.96 & 0.51 & 0.58 & 0.116 & 9.6\\
USMV replica, PCA & 0.98 & 1.09 & 0.56 & 0.68 & 0.109 & 8.7\\
USMV replica, beta exempt & 0.34 & 0.28 & -0.84 & -0.80 & 0.241 & 13.6\\
\addlinespace
SPLV, the fund & -1.19 & -0.56 & -3.35 & -1.95 & 0.472 & 14.0\\
SPLV replica & -0.48 & -0.24 & -2.50 & -1.51 & 0.448 & 14.0\\
\bottomrule
\end{tabular}
\notetext{Daily excess returns regressed on the Fama--French five factors plus momentum, then re-run with the betting-against-beta factor added. Annualised intercepts, with Newey--West standard errors at 8 lags. The two blocks end on different dates because the factor vintages do: the six-factor window runs 2011-05-06 to 2026-08-31 and the seven-factor window to 2026-07-31, so the columns are not a paired test on an identical sample. No intercept reaches significance at the five per cent level in either block, though the largest, the SPLV fund's own at $t=-1.95$, falls only just short at $p=0.051$ and would clear a ten per cent threshold. Every betting-against-beta loading is significant by a wide margin, which is why the paper reads the pattern as absorption rather than as an absence of exposure.}
\end{table}
